% !TeX root = ../../tesis.tex

\section{Proceso ETL: del GTFS a la Base de Datos}
\label{anx:apimetro-etl}

El proceso de ingesta opera en dos fases. La primera lee los archivos
\textsc{gtfs} originales y produce una plantilla Excel con los campos
normalizados al esquema de \texttt{Apimetro}. La segunda toma esa plantilla,
cruza las geometrías del archivo \texttt{stops.txt} y carga el resultado
directamente en PostgreSQL vía \textit{GeoPandas}.

\subsection*{Dato de entrada: \texttt{routes.txt} del \textsc{gtfs}}

El archivo \texttt{routes.txt} es uno de los cuatro archivos obligatorios del
estándar \textsc{gtfs}. Cada fila describe una línea de transporte con su
agencia, número comercial, nombre largo y colores institucionales. La
heterogeneidad de los valores en \texttt{agency\_id} —\textsc{semovi},
\textsc{metro}, \textsc{cbb}, \textsc{cc}, entre otros— es el origen de la
necesidad de normalización.

\lstinputlisting[ caption={Muestra del archivo \texttt{routes.txt} de la \textsc{adip} con nueve líneas de distintos sistemas.}, label=lst:gtfs-routes, basicstyle=\small\ttfamily] {Anexos/Apimetro-Anexos/ejemplo_GTFS_routes.txt}

\subsection*{Fase 1 — Generador de plantilla
(\texttt{ejemplo\_Generador\_Plantilla.py})}

El script lee los cuatro archivos \textsc{gtfs} requeridos, aplica la función de
estandarización de nombres por agencia y exporta una plantilla Excel con tres
hojas: \textit{Líneas}, \textit{Ramales} y \textit{Estaciones}. Se muestra la
función de normalización y el bloque de procesamiento de líneas; los bloques de
ramales y estaciones siguen el mismo patrón.

\lstinputlisting[language=Python, caption={Normalización de nombres de sistema a partir de \texttt{agency\_id} y generación de la plantilla Excel.}, label=lst:py-generador, breaklines=true, firstline=35, lastline=83, basicstyle=\small\ttfamily] {Anexos/Apimetro-Anexos/python/ejemplo_Generador_Plantilla.py}

\subsection*{Fase 2 — Carga a PostGIS
(\texttt{ejemplo\_DataEstacion.py})}

La clase \texttt{EstacionETL} implementa el patrón
\textit{Extract–Transform–Load} en cuatro métodos: lectura del Excel de
plantilla, lectura del \texttt{stops.txt}, cruce por \texttt{stop\_gtfs} y carga
final a la tabla \texttt{estacions} con \texttt{gdf.to\_postgis()}.

\lstinputlisting[language=Python, caption={Clase \texttt{EstacionETL}: extracción, transformación y carga a PostGIS vía \textit{GeoPandas}.}, label=lst:py-etl, breaklines=true, basicstyle=\small\ttfamily] {Anexos/Apimetro-Anexos/python/ejemplo_DataEstacion.py}
